\documentclass[12pt,a4paper]{article}
\usepackage{amsmath,amssymb}
\usepackage{booktabs}
\usepackage{geometry}
\usepackage{graphicx}
\usepackage{float}
\usepackage{listings}
\usepackage{xcolor}
\usepackage{algorithm}
\usepackage{algpseudocode}
\usepackage{titlesec}
\usepackage{setspace}
\usepackage{ragged2e}

%----------------------------------------
% Page Setup for A4
%----------------------------------------
\geometry{
  a4paper,
  top=1in,
  bottom=1in,
  left=1in,
  right=1in
}

\setstretch{1.2}

%----------------------------------------
% Section title formatting
%----------------------------------------
\titleformat{\section}
  {\normalfont\bfseries\fontsize{14}{14}\selectfont}
  {\thesection}{1em}{}

%----------------------------------------
% Code listing style
%----------------------------------------
\lstdefinestyle{matlabstyle}{
  language=Matlab,
  basicstyle=\ttfamily\footnotesize,
  keywordstyle=\color{blue}\bfseries,
  commentstyle=\color{green!60!black},
  stringstyle=\color{orange!80!black},
  frame=none,
  numbers=none,
  breaklines=true,
  showstringspaces=false,
  tabsize=2
}

%----------------------------------------
\begin{document}

%----------------------------------------
% Cover Page
%----------------------------------------
\begin{titlepage}
\centering

\vspace*{0.5cm}
\includegraphics[width=0.22\textwidth]{diu_logo.png}\par
\vspace{0.5cm}

{\Large \textbf{Department of Computer Science and Engineering} \par}
\vspace{0.2cm}
{\Large \textbf{Dhaka International University} \par}

\vspace{0.8cm}

{\normalsize
Semester: 12th \par
Lab Report No: 06 \par
Course Code: CSE-420 \par
Course Title: Image Processing Lab \par
Report Title: Edge Detection, K-Means Clustering, Connected Component Labeling, and Thresholding in MATLAB \par
}

\vspace{1cm}

{\Large \textbf{Submitted By} \par}
\vspace{0.3cm}

{\normalsize
Rasheduzzaman Rakib \par
Reg. No.: CS-D-77-22-120080 \par
Roll: 37 \par
Batch: D-77 \par
Department of Computer Science and Engineering \par
Dhaka International University \par
}

\vspace{1cm}

{\Large \textbf{Submitted To} \par}
\vspace{0.3cm}

{\normalsize
Md. Muksit Ul Islam \par
Assistant Professor \par
Department of Computer Science and Engineering \par
Dhaka International University \par
}

\vfill

{\normalsize Submission Date: 25th September 2026 \par}
{\normalsize Summer 2026 \par}

\end{titlepage}


%----------------------------------------
% Lab Report
%----------------------------------------

\section{Name}

\textbf{Edge Detection, K-Means Clustering, Connected Component Labeling, and Thresholding in MATLAB}\\[0.3em]
Topics covered: Canny Edge Detection, Prewitt Edge Detection, K-Means Color Segmentation, Object Extraction, Manual Thresholding, Connected Component Labeling, and Otsu Thresholding.

%----------------------------------------

\section{Description}

This lab demonstrates edge detection, unsupervised pixel clustering, and object segmentation methods in MATLAB using peppers.png, cameraman.tif, and coins.png.

\begin{enumerate}
  \item \textbf{Canny Edge Detection:} Identifies prominent edges through Gaussian smoothing, directional gradient computation, non-maximum suppression, and hysteresis thresholding via edge with canny parameter.

  \item \textbf{Prewitt Edge Detection:} Uses 3$\times$3 horizontal and vertical gradient operator masks to locate intensity boundaries in grayscale images.

  \item \textbf{K-Means Clustering:} Groups RGB pixel vectors into $K = 4$ clusters using the kmeans algorithm to segment regions having similar color profiles.

  \item \textbf{Object Extraction and Color Quantization:} Creates a binary cluster mask to isolate peppers against a black background and remaps pixels to their centroid colors.

  \item \textbf{Manual Thresholding:} Binarizes the cameraman image by applying a fixed user-defined cutoff value (0.2) via im2bw.

  \item \textbf{Connected Component Labeling:} Applies bwlabel on thresholded coins image to label contiguous pixel groups and return the total number of connected objects.

  \item \textbf{Otsu Automatic Thresholding:} Computes the optimal threshold with graythresh to cleanly separate foreground objects from the background automatically.
\end{enumerate}

%----------------------------------------

\section{Code}

\begin{lstlisting}[style=matlabstyle]
% read images
I = imread('peppers.png');
I1 = rgb2gray(I);
C = imread('cameraman.tif');
coins = imread('coins.png');

% 1. canny edge detection
BW1 = edge(I1, 'canny');

% 2. prewitt edge detection
BW2 = edge(I1, 'prewitt');

% 3. k-means clustering
X = double(reshape(I, [], 3));
K = 4;
[idx, C1] = kmeans(X, K);
cluster_img = reshape(idx, size(I,1), size(I,2));
pepper_cluster = 2;
mask = (cluster_img == pepper_cluster);
result = zeros(size(I), 'uint8');

for k = 1:3
    temp = I(:,:,k);
    temp(~mask) = 0;
    result(:,:,k) = temp;
end

% 4. color k-means
new_image = C1(idx,:);
new_image = uint8(reshape(new_image, size(I)));

figure;
subplot(2,3,1);
imshow(I);
title('original image');

subplot(2,3,2);
imshow(BW1);
title('canny edge detection');

subplot(2,3,3);
imshow(BW2);
title('prewitt edge detection');

subplot(2,3,4);
imshow(cluster_img,[]);
title('k-means clusters');

subplot(2,3,5);
imshow(result);
title('peppers with black background');

subplot(2,3,6);
imshow(new_image);
title('color k-means segmentation');

% 5. manual thresholding
BW3 = im2bw(C, 0.2);

% 6. connected component labeling
BW4 = im2bw(coins, graythresh(coins));
[L,num] = bwlabel(BW4);

% 7. otsu automatic thresholding
level = graythresh(C);
BW5 = im2bw(C, level);

% display remaining results
figure;
subplot(2,3,1);
imshow(C);
title('original image');

subplot(2,3,2);
imshow(BW3);
title('manual threshold');

subplot(2,3,3);
imshow(coins);
title('coins image');

subplot(2,3,4);
imshow(L,[]);
title('connected components');

subplot(2,3,5);
imshow(C);
title('original image');

subplot(2,3,6);
imshow(BW5);
title('otsu segmentation');

disp('number of connected components =');
disp(num);
\end{lstlisting}

%----------------------------------------

\section{Output}

Figures~\ref{fig:edges_kmeans} and~\ref{fig:thresh_cc} present the output figures obtained from the execution of the code.

\begin{figure}[H]
  \centering
  \includegraphics[width=0.7\textwidth]{Figure_1.png}
  \caption{Edge detection and K-means color clustering on peppers.png: original, Canny edges, Prewitt edges, cluster index map, isolated pepper object, and color-quantized segmentation.}
  \label{fig:edges_kmeans}
\end{figure}

\begin{figure}[H]
  \centering
  \includegraphics[width=0.7\textwidth]{Figure_2.png}
  \caption{Segmentation and connected components: manual threshold, connected component labeling on coins.png, and Otsu automatic thresholding on cameraman.tif.}
  \label{fig:thresh_cc}
\end{figure}

Figure~\ref{fig:edges_kmeans} demonstrates that Canny detects fine, thin, continuous boundaries superior to Prewitt, while K-means cleanly clusters chromatic features. Figure~\ref{fig:thresh_cc} shows bwlabel differentiating and counting each individual coin component, whereas Otsu thresholding delivers superior separation compared to manual thresholding.

\end{document}
